% =====================================================
% 题型：正方体线线夹角选择题 (q_solid_cube_plane_intersection_angle.tex)
% =====================================================
% =====================================================
% 难度：★★ (中档)
% =====================================================
\newcommand{\qSolidCubePlaneIntersectionAngleStem}{如图，在正方体 \(ABCD-A_1B_1C_1D_1\) 中，过直线 \(AA_1\) 且平行于平面 \(BDD_1B_1\) 的平面 \(\alpha\) 交平面 \(ACD_1\) 于直线 \(l\)，则直线 \(AA_1\) 与 \(l\) 所成角的正切值为 \hfill （\quad）}

\newcommand{\qSolidCubePlaneIntersectionAngleOptions}{%
  \mcoptions[4]{\(\dfrac12\)}{\(\dfrac{\sqrt2}{2}\)}{\(\sqrt2\)}{\(2\)}%
}

\newcommand{\qSolidCubePlaneIntersectionAngleAnswer}{B}

\newcommand{\qSolidCubePlaneIntersectionAngleFigure}[1][0.85]{%
  \begin{tikzpicture}[scale=#1, line join=round, line cap=round, >=stealth]
    % Coordinates of vertices (oblique projection)
    % A(0,0), B(2.2,0), C(3.0,0.8), D(0.8,0.8)
    \coordinate (A) at (0,0);
    \coordinate (B) at (2.2,0);
    \coordinate (C) at (3.0,0.8);
    \coordinate (D) at (0.8,0.8);
    \coordinate (A1) at (0,2.2);
    \coordinate (B1) at (2.2,2.2);
    \coordinate (C1) at (3.0,3.0);
    \coordinate (D1) at (0.8,3.0);
    
    % Back vertices and lines (dashed)
    \draw[thick, dashed] (D) -- (A);
    \draw[thick, dashed] (D) -- (C);
    \draw[thick, dashed] (D) -- (D1);
    \draw[thick, dashed] (B) -- (D) node[midway, above left, xshift=2pt] {};
    \draw[thick, dashed] (B1) -- (D1);
    
    % Front vertices and lines
    \draw[thick] (A) -- (B) -- (C) -- (C1) -- (B1) -- (A1) -- cycle;
    \draw[thick] (A1) -- (B1) -- (C1);
    \draw[thick] (B) -- (B1);
    
    % plane ACD1
    \draw[thick, dashed, gray] (A) -- (C);
    \draw[thick, dashed, gray] (A) -- (D1);
    \draw[thick, dashed, gray] (C) -- (D1);
    
    % Line AA1
    \draw[very thick, double] (A) -- (A1);
    
    % Labels
    \node[below left] at (A) {\small \(A\)};
    \node[below right] at (B) {\small \(B\)};
    \node[right] at (C) {\small \(C\)};
    \node[above left] at (D) {\small \(D\)};
    \node[above left] at (A1) {\small \(A_1\)};
    \node[above right] at (B1) {\small \(B_1\)};
    \node[above right] at (C1) {\small \(C_1\)};
    \node[above] at (D1) {\small \(D_1\)};
    
    \ifteacher
      % Show plane alpha and line l
      % plane alpha is x+y=0.
      % intersection line l is (1, -1, -2) passing through A.
      % To draw it in oblique projection:
      % direction (1, -1, -2) in A coordinate system is (B - D - 2*A1)
      % x-direction is parallel to AB (2.2, 0)
      % y-direction is parallel to AD (0.8, 0.8)
      % z-direction is parallel to AA1 (0, 2.2)
      % So vector is (2.2 - 0.8, 0 - 0.8, -4.4) = (1.4, -0.8, -4.4)
      % This goes downwards, let's draw it from A.
      % To make it look like a line, we can draw the line l starting from A(0,0) and going to (-0.7, 0.4, 2.2) = (-0.7 - 0.4, 0.4 - 0.4, 2.2) = (-1.1, 0, 2.2)
      % In our 2D coordinates, let's plot l direction: direction is (-1.1, 2.2) relative to A.
      \coordinate (L_dir) at (-0.8, 1.6);
      \draw[very thick, red, ->] (A) -- (L_dir) node[left] {\small \(l\)};
      \draw[thick, dashed, blue] (A) -- (-0.8, 0.0) node[below] {\small 平面 \(\alpha\) 迹线};
    \fi
  \end{tikzpicture}%
}

\newcommand{\qSolidCubePlaneIntersectionAngleSolution}{%
  设正方体棱长为 \(1\)，以 \(A\) 为原点，\(AB, AD, AA_1\) 所在直线分别为 \(x, y, z\) 轴建立空间直角坐标系。
  
  则各点坐标为：
  \[
    A(0,0,0),\quad B(1,0,0),\quad C(1,1,0),\quad D(0,1,0),\quad A_1(0,0,1),\quad D_1(0,1,1).
  \]
  
  平面 \(BDD_1B_1\) 经过点 \(B(1,0,0)\), \(D(0,1,0)\)，且平行于 \(z\) 轴（\(AA_1\) 方向），其方程为：
  \[
    x + y = 1.
  \]
  
  平面 \(\alpha\) 过直线 \(AA_1\) 且平行于平面 \(BDD_1B_1\)。
  因为直线 \(AA_1\) 在 \(z\) 轴上（\(x=0, y=0\)），且 \(\alpha \parallel \text{平面 } BDD_1B_1\)，所以平面 \(\alpha\) 的方程为：
  \[
    x + y = 0.
  \]
  
  接着求平面 \(ACD_1\) 的方程。因为 \(A(0,0,0)\)，所以平面过原点。
  向量 \(\overrightarrow{AC} = (1,1,0)\)，\(\overrightarrow{AD_1} = (0,1,1)\)。
  设平面 \(ACD_1\) 的法向量为 \(\vec{n} = (x_0, y_0, z_0)\)，则：
  \[
    \begin{cases}
      \vec{n} \cdot \overrightarrow{AC} = x_0 + y_0 = 0 \\
      \vec{n} \cdot \overrightarrow{AD_1} = y_0 + z_0 = 0
    \end{cases}
    \implies \vec{n} = (1, -1, 1).
  \]
  所以平面 \(ACD_1\) 的方程为：
  \[
    x - y + z = 0.
  \]
  
  直线 \(l\) 为平面 \(\alpha\) 与平面 \(ACD_1\) 的交线，联立两平面方程：
  \[
    \begin{cases}
      x + y = 0 \\
      x - y + z = 0
    \end{cases}
    \implies y = -x, \quad z = -2x.
  \]
  取 \(x = 1\)，则 \(y = -1\), \(z = -2\)。
  所以直线 \(l\) 的一个方向向量为：
  \[
    \vec{l} = (1, -1, -2).
  \]
  
  而直线 \(AA_1\) 的方向向量为：
  \[
    \vec{a} = (0, 0, 1).
  \]
  
  设直线 \(AA_1\) 与 \(l\) 所成的角为 \(\theta\)，由向量夹角公式：
  \[
    \cos\theta = \frac{|\vec{l} \cdot \vec{a}|}{|\vec{l}||\vec{a}|} = \frac{|-2|}{\sqrt{1^2 + (-1)^2 + (-2)^2} \cdot \sqrt{1}} = \frac{2}{\sqrt{6}} = \sqrt{\frac{2}{3}}.
  \]
  
  因为 \(\theta \in \left[0, \dfrac{\pi}{2}\right]\)，所以：
  \[
    \sin\theta = \sqrt{1 - \cos^2\theta} = \sqrt{1 - \frac{2}{3}} = \frac{1}{\sqrt{3}}.
  \]
  
  因此，两直线所成角的正切值为：
  \[
    \tan\theta = \frac{\sin\theta}{\cos\theta} = \frac{1/\sqrt{3}}{\sqrt{2}/\sqrt{3}} = \frac{1}{\sqrt{2}} = \frac{\sqrt{2}}{2}.
  \]
  故选 B。%
}

% =====================================================
% 别名兼容：旧课次及学生个人宏定义
% =====================================================
\newcommand{\qTJWCubePlaneIntersectionAngleStem}{\qSolidCubePlaneIntersectionAngleStem}
\newcommand{\qTJWCubePlaneIntersectionAngleOptions}{\qSolidCubePlaneIntersectionAngleOptions}
\newcommand{\qTJWCubePlaneIntersectionAngleFigure}{\qSolidCubePlaneIntersectionAngleFigure}
\newcommand{\qTJWCubePlaneIntersectionAngleAnswer}{\qSolidCubePlaneIntersectionAngleAnswer}
\newcommand{\qTJWCubePlaneIntersectionAngleSolution}{\qSolidCubePlaneIntersectionAngleSolution}
